The army had stopped looking like an army somewhere around the third day north of White Harbor, and had started looking, instead, like a long wound dragging itself across white country, and Jon Snow, riding at its head with his hood drawn low against a wind that seemed to find its way beneath every layer of wool and leather a man owned, had begun to wonder whether that was not, in truth, the more honest shape for it to take.

The Unsullied bore the cold the way they bore everything — in silence, in unbroken ranks, spears shouldered at the same precise angle whether the sun stood high or the snow came sideways hard enough to blind a man three feet in front of him. Grey Worm walked them in shifts through the worst of the drifts, rotating the outer files inward before frostbite could claim toes no soldier would think to complain of losing, and Jon had come, these past weeks, to a grim respect for that particular kind of discipline — not the absence of suffering, he understood now, but the refusal to let suffering announce itself.

The Dothraki suffered where the Unsullied merely endured, and did not trouble to hide it. Jon had ridden alongside Jorah Mormont two mornings past to find the old bear of Bear Island doing what Jon suspected no other man in Daenerys’s service could have done — moving horse to horse among Khal Drogo’s bloodriders, speaking low in a Dothraki roughened by years but still serviceable, showing them how to pack snow from a horse’s hooves before it froze there and split the frog, how to blanket a mount properly against wind that would otherwise kill it by morning. It was strange to watch, Jon thought — a Northman’s knowledge, learned as a boy on Bear Island’s own hard winters, spent now keeping horse-lords from the Dothraki Sea alive in a cold their own gods had never troubled to warn them about.

“They lost eleven horses last night,” Jorah told him, when Jon asked, his breath fogging thick in air cold enough to hurt the lungs that breathed it. “Would have lost thirty, the pace we were losing them a week back. They’re proud people, Lord Snow. Prouder than sense allows, some of them, and prouder still than the North permits any man to be who wants to see spring again. But they’re learning. Slower than I’d like. Faster than I feared.”

“And you? How does Bear Island’s own son find his homecoming?”

Something crossed the old knight’s weathered face that Jon couldn’t quite read — not warmth, not quite grief either. “Fifteen years an exile, coming home to a war that may kill every soul I left behind twice over, in a season that never seems to end.” A short, humorless breath that might have been meant as a laugh. “I find it colder than I remembered, Lord Snow, in more ways than the one.”

Jon did not ask him to say more. There was a look men got, coming home to things they were no longer certain still wanted them, and he had worn that look himself often enough to know better than to make another man explain it aloud.

\sceneBreak

Davos found him not long after, picking his way across ground already churned to black mud and grey slush by eight thousand new feet, a look on his face that Jon had learned, across two years, to read as trouble before the man had said a word of it.

“Lord Manderly’s steward wants a word,” Davos said. “Politely put, it’s a question of grain. Less politely put, it’s a Northman standing in a granary door asking why his lord’s stores should feed Dothraki horses before they feed Northern children, and I didn’t have a good enough answer for him that didn’t sound like something a Southron lord says right before he takes what he wants regardless.”

“What did you tell him?”

“I told him we’d discuss it properly, in the hall, where the whole matter could be weighed fairly instead of settled by one steward and one Onion Knight arguing in a doorway.” Davos’s mouth twisted, not quite a smile. “He didn’t much care for that answer either. I don’t imagine he’s alone in it. There’s forty thousand new mouths in this castle and the country around it tonight, Jon, counting horses, and not one of them Northern-born. Men notice arithmetic like that faster than they notice politics, and they’ll forgive politics a good deal longer than they’ll forgive going hungry for it.”

“Then we’d best have a real answer ready, before it’s a steward asking and not a lord.”

“Aye,” Davos said, grim. “I expect we will, and soon.”

\sceneBreak

He found Grey Worm and Missandei that same evening, at the very edge of the day’s march where the cookfires hadn’t yet reached to crowd them, and slowed his own steps rather than intrude on what was plainly meant to be a private moment — Missandei wrapped in furs three sizes too large for her, laughing at something Grey Worm had said too low for Jon to catch, her breath fogging white against a sky gone the deep bruised purple of a Northern dusk. There was an ease between them that the cold and the war and the whole grinding weight of the march hadn’t managed to touch, and Jon found himself, watching it a moment longer than he’d meant to, thinking of a boat’s cabin some weeks back and a laugh that had caught him as much by surprise as it seemed to have caught the woman who gave it.

He moved on before either of them noticed him watching, because some things a man let people keep for themselves, war or no war, and because he had his own reunions still waiting for him at the end of this day’s ride, and found, walking back toward the column’s head, that he was more afraid of those than he had been of anything the wolfswood’s dark had offered on the ride north.

\sceneBreak

Winterfell’s walls came up out of the last grey light like something out of an old memory sharpened too suddenly into focus, and Jon felt the whole column behind him go quiet in a way that had nothing to do with discipline and everything to do with awe, because Rhaegal and the silver queen’s shadow overhead had chosen that moment to bank low over the battlements, wings catching the day’s last light, and every face along Winterfell’s walls that Jon could make out from this distance had turned upward, pale and open-mouthed, smallfolk and soldiers alike struck by the same old terror the North’s stories had always promised and none of them, until this moment, had ever quite believed.

Arya was waiting inside the gate. He knew her before he could properly see her face — the way she stood, weight balanced forward on the balls of her feet even at rest, a habit that had not belonged to the girl who left King’s Landing all those years ago and had plainly been built into whoever she’d become since — and when she saw him she did not run to him the way the girl he remembered might have, but walked, unhurried, until she was close enough that the last few steps became, almost against her own evident will, something closer to running after all.

“You got taller,” she said, into his shoulder, muffled, and he laughed, the sound surprising him with how much he’d needed it.

“You got faster. And quieter. I didn’t hear you at all, coming up the yard.”

“You weren’t meant to.” She pulled back to look at him properly, and he saw, in the assessment behind her eyes, something that was not quite suspicion and not quite welcome either — an accounting, the kind he imagined she’d learned to take of every stranger she met these past years, applied now, a little helplessly, even to her own brother. “You brought a queen. And dragons. And an army that isn’t ours.”

“I brought what I could get, to fight what’s coming. I’d rather have brought nothing and had it be enough. It wasn’t.”

“No,” Arya agreed, quiet, some of the guardedness easing though not all of it. “I don’t imagine it ever is.” Her hand found the hilt at her hip, unconsciously, the way Jon remembered her doing as a girl with a much smaller blade. “Needle’s held up. Thought you’d want to know that, since you’re the one who gave it to me.”

“I’d have expected nothing less.”

“You should. I’ve killed men who expected less.” She said it flatly, not for shock, simply as fact, the way she might have once said she’d caught a fish or won a footrace, and Jon felt the small, cold weight of everything he didn’t yet know about the years she’d spent away from him settle somewhere behind his ribs, unexamined, waiting for a quieter hour that the war might or might not permit them.

He left her at the practice yard a while later, still turning Needle in her hand as though reacquainting herself with an old friend’s grip, and found her, not an hour after, in close conversation with Lyanna Mormont — the Lady of Bear Island barely past her tenth year in body, and old as any soldier Jon had ever commanded in the flat, unblinking way she held a room. The two of them made an odd pair, seen from a little distance — one a woman grown hard and quiet in foreign cities Jon had never set foot in, the other a girl who had never once in her short life been permitted the luxury of being only a girl — and yet something in the set of their shoulders matched so precisely that Jon slowed to watch rather than interrupt.

“They tell songs about you already,” Lyanna was saying, without any of the awe a girl her age might have been expected to bring to it. “Killed a Faceless Man’s teacher and wore his face after. Killed a Frey lord at his own wedding table and served him to his own kin in a pie.”

“Songs get details wrong,” Arya said. “Usually the interesting ones.”

“Which details did they get wrong?”

“That’s between me and the men who deserved it.” But there was something almost like respect in the look Arya gave her, one soldier’s flat assessment of another, regardless of years. “You held Bear Island’s men at Winterfell’s first battle. Stood in a shield line at an age most girls are still afraid of the dark.”

“I was never afraid of the dark,” Lyanna said, with the particular contempt of someone who had genuinely never had the luxury. “I was afraid of what my house would become if I let anyone see fear on me first. Turns out that’s most of what commanding men actually is.”

“Turns out,” Arya agreed, “that’s most of what staying alive actually is, too.” And for a moment, watching them, Jon thought the two of them understood something about each other — about the particular cost of becoming, too young, exactly hard enough to survive what the world had asked of them — that neither Sansa’s careful diplomacy nor his own soldier’s discipline had ever quite required of either of them in the same unforgiving way.

He left them to it, and did not interrupt.

\sceneBreak

Samwell Tarly found him near the armory, arms full of books he seemed unable to set down even to embrace an old friend properly, and the sight of him — heavier than the boy who’d fled the Wall with him a lifetime ago, but standing straighter too, in a way Jon suspected had less to do with his body than with whatever he’d been reading these past months — startled a real laugh out of Jon for the first time since White Harbor.

“You’re meant to be a maester by now, and you still can’t carry your books without dropping half of them.”

“I dropped exactly one, and only because you startled me.” Sam’s grin faded some as he set the books down properly at last, and Jon saw, beneath the familiar softness of his old friend’s face, something new and harder settling in around the eyes. “Gilly and little Sam are safe, if that’s the first thing you’re going to ask, which I know it is. Inside the walls, warm, fed. I’ve seen to it myself.”

“And you? What have you found, in all those books, that’s put that look on your face?”

Sam hesitated, glancing toward the godswood as though the answer lived there rather than in his own throat. “Fragments, mostly. Half-legends the maesters never troubled to copy twice, because no one thought they’d matter again. Things about the Long Night that don’t quite fit together the way a proper history should. I keep thinking Bran could make more sense of it than I can, if I could ever get him to sit still with me long enough to try, but he’s—“ Sam stopped himself, seemed to reconsider how much he wanted to say plainly, out here in the open yard. “He’s different, Jon. You’ll see it yourself, soon enough.”

“I’ve heard as much already.”

“Hearing it and seeing it aren’t the same.” Sam’s hand found Jon’s shoulder, brief, warm, entirely unchanged from the boy who’d sworn brotherhood with him at the Wall a hundred years and one crippled kingdom ago. “I’m glad you’re back all the same. Whatever else is coming, I’m glad of that much.”

“So am I,” Jon said, and meant it, and left Sam to his books and went on toward the godswood, where the last of his reunions still waited, colder and stranger than any of the ones that had come before it.

\sceneBreak

Bran received him in the godswood, as he seemed now to receive everyone, and Jon, crossing that familiar ground toward a brother he had not laid eyes on since a summer that felt, now, like it belonged to some other man’s life entirely, found himself slowing without quite meaning to, the way a man slows approaching a door he isn’t certain he wants opened.

Bran did not rise, because Bran could not rise, and that alone struck Jon harder than he’d braced for — the wheeled chair, the furs, the stillness of legs that would never again carry a boy up Winterfell’s walls the way they once had. But it was the eyes that unsettled him past bracing, once he was close enough to see them properly: patient in a way no boy of Bran’s years had any right to be, focused somewhere just past Jon’s shoulder as often as on his face, as though half of whatever sat in that chair were still listening to some other conversation entirely, happening somewhere Jon could not follow.

“Bran.” He knelt, so they were level, because it seemed wrong, suddenly, to stand over him. “Gods, it’s good to see you. I never thought—“ He stopped, because he wasn’t certain what he’d never thought, or which of a dozen half-formed dreads he meant to name.

“You never thought I’d live,” Bran said, without unkindness, “or you never thought, if I lived, that I’d still be the boy who fell from that tower. I’m neither of the things you’re afraid to say, Jon. I’m something else now. I’m trying to find better words for what, but the words keep arriving slower than what I actually am.”

“You knew I was coming. Before I sent word.”

“I know a great many things before anyone sends word of them. It isn’t as useful as it sounds, most days. I can tell you the shape of the storm and still not stop it from breaking.” Bran’s eyes found his properly then, held them, and Jon felt, absurdly, like a much younger man being seen through entirely. “You’re afraid of the throne you gave up. Not for yourself. For what happens after, if the woman you knelt to turns out to want a different ending than the one you’re hoping for.”

Jon did not answer that. There did not seem, in that moment, to be an answer worth giving that Bran had not already read off him before he’d thought to hide it, and he knelt there a while longer in the cold godswood, brother across from brother, saying nothing at all.